\section{Corpus and Methodology}
\label{sec:method}

\subsection{Reproducible environment}

All results are real runs in one pinned container: Ubuntu~24.04 on
\texttt{amd64}, GCC~13.3.0, Clang~18.1.3, Valgrind (Memcheck). A verdict can
depend on compiler and optimization level, so one without the exact compiler
version is under-specified. The suite has 381 passing and 3 skipped tests; skips
need a native (non-emulated) timing environment.

\subsection{Corpus design}

The corpus pairs each layer with a case \emph{only} that layer flags, grounding
every verdict class: real PQClean targets~\cite{pqclean} plus synthetic
controls (Table~\ref{tab:coverage}, per-feature map;
Table~\ref{tab:corpus}, verdict-class corpus).

\paragraph{Real PQC targets.}
For ML-KEM-768 (FIPS~203~\cite{fips203}, \texttt{mlkem768}, harness
\texttt{kem\_dec}) we screen the plain build and a
KyberSlash~\cite{kyberslash} variable-time-division variant
(\texttt{mlkem768\_kyberslash}). The plain build passes structurally
everywhere; \texttt{asm-scan}'s FIPS~202/Keccak division candidates triage
\emph{public} (constant absorb rate). KyberSlash is the key case: Memcheck
passes it clean---no secret-tainted branch or address---yet it leaks through a
secret-derived integer division only \texttt{asm-scan} surfaces.

\paragraph{Synthetic controls.}
Three harnesses isolate one phenomenon each. \texttt{toy\_lookup}'s
\texttt{leaky} variant (\texttt{out[i]=sbox[secret[i]]}) is a cache-timing leak
Memcheck flags across all six builds; \texttt{safe} is robust.
\texttt{ct\_matrix\_flip}'s secret-dependent branch is optimizer-removed: it
fails under GCC debug (\texttt{-O0}) but passes under release (\texttt{-O2}) and
size (\texttt{-Os}). \texttt{toy\_dudect} is a timing control, no structural
branch (Table~\ref{tab:dudect}).

\subsection{How a row is produced}

Each row follows one pipeline (Fig.~\ref{fig:pipeline}) from one YAML config.
\texttt{ct-matrix} recompiles per cell of the compiler~$\times$ named-cflags
matrix (up to six cells, GCC and Clang) using the debug/release/size
configurations of Section~\ref{ssec:matrix}, re-running the \emph{same} check in
each. \texttt{asm-scan} runs warn-only over object code; public
variable-latency candidates are recorded, any unrecognized leak site holds the
row at \texttt{needs-analysis}, and the default-deny taxonomy never auto-promotes
a structural failure to a leak. Optionally \texttt{dudect}~\cite{dudect} (a Welch
$t$-test on cycle counts) runs on timing distributions. The ML-DSA-65 row
(\texttt{mldsa65}, \texttt{sign}) is the central honest result: its structural
check fails, but default-deny holds it at \texttt{needs-analysis} rather than
declaring a leak (Section~\ref{sec:results-coverage}).

\subsection{Merging the evidence}

Structural and \texttt{asm-scan} signals are per build cell, the timing test per
target. A script merges them, collapsing structural results into one verdict per
target (build-sensitive when verdicts vary) and attaching triaged findings:
Table~\ref{tab:corpus}, seven rows, three families, five verdict classes. The timing layer is in the dudect appendix (Table~\ref{tab:dudect}).
